\chapter{Récoltes}
\label{chap:recoltes}

\section{Récoltes annuelles}

Pour les plantations \textbf{annuelles}, la fenêtre de récolte est définie au moment de la création (date de début et date de fin de récolte). Ces jalons apparaissent dans le \textbf{Calendrier} unifié (en contour) et une tâche \emph{Récolte} est générée à la date de début (voir chapitre \ref{chap:calendrier}).

\begin{notebox}
\textbf{Prévu.} La saisie d'événements de récolte ponctuels (date, quantité pesée) pour les plantations annuelles, à la manière de Qrop, est prévue pour une version ultérieure.
\end{notebox}

\section{Récoltes annuelles pour plantations pluriannuelles}

Pour les plantations \textbf{pluriannuelles} (vergers, petits fruits, vivaces), chaque année donne lieu à une entrée de récolte distincte. Depuis la vue détail d'une plantation pluriannuelle, vous pouvez enregistrer :

\begin{itemize}
  \item L'\textbf{année} de récolte.
  \item Le \textbf{rendement attendu} (kg) pour cette année.
  \item Le \textbf{rendement réalisé} (kg).
  \item Des \textbf{notes} libres.
\end{itemize}

Pomone calcule automatiquement l'écart estimé/réalisé. L'historique des années est conservé pour analyse pluriannuelle ; ré-enregistrer une année existante met l'entrée à jour.

% TODO capture — pré-prod : ajouter une capture du tableau des récoltes annuelles
% et de son formulaire d'ajout (données réelles).
